\section{Liquidated damages strategy}

\subsection{Design principles}
\begin{enumerate}\setlength{\itemsep}{0.3em}
  \item \textbf{Express LDs as a percentage of OEM supply value, not absolute
        dollars.} OEMs assess exposure relative to project value; a percentage
        basis is more persuasive and scales correctly.
  \item \textbf{Set LDs high enough to change behavior.} If the penalty is too
        low it becomes more convenient for the OEM to pay out and deprioritize
        American Terawatt's project in favor of higher-penalty contracts. The
        purpose of the LD is to \emph{incentivize the OEM to fix problems}.
  \item \textbf{Do not expect full recovery.} LDs will not recover the true
        commercial loss from unavailability; they are a behavioral lever, and
        the residual exposure must be managed elsewhere (redundancy, storage).
\end{enumerate}

\subsection{Recommended approach}
\begin{itemize}\setlength{\itemsep}{0.25em}
  \item \textbf{Anchor the LD basis to the PPA.} Research what the data center
        charges American Terawatt for unavailability under the PPA, and use a
        fair share of that as the LD basis --- a concrete, defensible argument
        to the OEM.
  \item \textbf{Quantify with battery cost.} Frame the FEU LD around the cost of
        the additional battery capacity needed to compensate for forced
        energy-unavailability events; this gives a strong, quantifiable figure.
  \item \textbf{Argue for a higher cap.} Rather than accepting the standard
        3--5\,\% cap, argue for a cap tied to the overall PPA liability that
        American Terawatt itself carries.
\end{itemize}

\subsection{Residual exposure}
Even with a well-designed LD regime, American Terawatt retains exposure to the
difference between recovered LDs and actual loss. Mitigation levers include
schedule and spares strategy, redundancy in the transmission configuration, and
battery/on-site generation sizing.
